\begin{claim}{C.51}{Proved}{Cocompact proper operators control the entire frontier}
Let $f_t$ be the based ambient-null plane loop of C.49,
let $G=\langle f_t\rangle$, and let $H$ commute with all
its values. For an open $G$-orbit $O$, put
\[
 H_O=\{h\in H:hO=O\},\qquad
 Q_O=\overline{\{h|_O:h\in H_O\}}^{\operatorname{Homeo}(O)}.
                                                        \tag{C.51.1}
\]
Consider all such $O$ for which $O/Q_O$ is compact and
the full family $\mathscr F_O$ of actual lifted boundary
quotient fiber partitions in C.46 has at most two members.
For every one of these orbits the restricted loop has
a based contraction on its entire closure,
\[
 \overline K^O_{s,t}\in\operatorname{Homeo}(\overline O),
 \qquad
 \overline K^O_{0,t}=f_t|_{\overline O},\qquad
 \overline K^O_{1,t}=\overline K^O_{s,0}
                         =\overline K^O_{s,1}=\id,
                                                        \tag{C.51.2}
\]
jointly continuous with its inverse. It commutes with
$H_O$ and equals exactly the C.34 median contraction
on every nonopen orbit in $\partial O$.
When $\mathscr F_O$ is empty, every point of
$\partial O$ is fixed by $G$.

The contractions can be chosen with exact transport
under $H$. They agree on all common frontiers and with
C.34, and hence paste over the closed nonopen locus
together with any finite collection of these closures.
This assertion does not prescribe their values on the
closed set $X_*$ of C.50: the comparison there with
its exponential prescription remains to be proved.
\end{claim}

\begin{proof}
Write $Q=Q_O$ and $L=\overline G$ in the plane
homeomorphism group. C.30 proves that the actual group
$Q$ in (C.51.1) is closed, locally compact, free and
proper on $O$, with the specified compact-open topology.
The restriction kernel of $H_O\to Q$ acts trivially
on $O$ and therefore on its closure.
We retain that map; we do not replace an element of
$Q$ by an unspecified global extension.

\emph{A compact set meeting every operator orbit.}
The quotient map $O\to O/Q$ is open. For each point
of $O$, choose a relatively compact open neighborhood
inside $O$. The images of these neighborhoods cover
the compact quotient. Choose finitely many whose
images still cover it, and let $C\subset O$ be the
union of their closures. Then $C$ is compact and
\[
                              Q C=O.                   \tag{C.51.3}
\]
This argument supplies the exact compact fundamental
set required below; it uses no fundamental-domain
boundary or continuous choice of orbit representatives.

\emph{Every frontier sequence produces an actual quotient.}
Take a sequence $x_n\in O$ converging in the plane to
$z\in\partial O$. By (C.51.3), write
$x_n=q_ny_n$ with $q_n\in Q$ and $y_n\in C$.
Pass to a subsequence with $y_n\to y\in C$.
For the moment choose $h_n\in H_O$ with
\[
                    |h_ny_n-q_ny_n|<1/n.               \tag{C.51.4}
\]
Such a choice follows from the definition of $Q$ as
the closure of the actual restrictions in (C.51.1).
Here and below $|\cdot|$ is the Euclidean distance
in the original ambient plane.

Every $h_n$ commutes with $L$: it commutes with all
generators of $G$, and commutation passes to its
closure in the plane homeomorphism group.
The orbit map
\[
                      L\longrightarrow O,\quad l\mapsto ly
                                                        \tag{C.51.5}
\]
is open by the same Polish orbit-map theorem used in
C.30 and C.42. Thus there are $l_n\to1$ in $L$ with
$l_ny=y_n$.
For clarity, take a decreasing countable identity
neighborhood basis $V_k$ in $L$. Each $V_k y$ is a
neighborhood of $y$, so $y_n\in V_k y$ eventually.
Choose an increasing sequence of such eventual
indices and then $l_n\in V_{k(n)}$ with
$k(n)\to\infty$. This proves the stated selection.
Commutation now gives the exact equality
\[
                       h_ny=l_n^{-1}h_ny_n\longrightarrow z.
                                                        \tag{C.51.6}
\]
Convergence follows from (C.51.4), $x_n\to z$, and
compact-open convergence of $l_n^{-1}$ to the identity
on a compact neighborhood of $z$.

Let $J_y=L_y$. For $j\in J_y$,
$j h_ny=h_nj y=h_ny$, so (C.51.6) gives $jz=z$.
Consequently $z\in\operatorname{Fix}(J_y)$, and
C.42's actual evaluation comparison gives
\[
 h_n|_O\longrightarrow\Phi_z^y
       \quad\hbox{uniformly on compact subsets of }O,
 \qquad
 \Phi_z^y:O\longrightarrow\R^2,\quad \Phi_z^y(ly)=lz.
                                                        \tag{C.51.7}
\]
This uses the homeomorphism
$\operatorname{Map}_L(O,\R^2)\to\operatorname{Fix}(J_y)$
of (C.42.1), including its compact-open topology.
In particular its assertion applies when $Gz$ is
a single point, in which case $\Phi_z^y$ is constant.
There is no convergence claim for arbitrary proper
operator maps outside their actual global approximants.

We compare (C.51.7) with the quotient family defined
at the fixed basepoint $x$ in C.46. Choose $a\in L$
with $ax=y$, and put $z_0=a^{-1}z$.
Then $z_0\in\operatorname{Fix}(L_x)\cap\partial O$.
For every $l\in L$ the equality
\[
             \Phi_{z_0}^x(ly)
                =\Phi_{z_0}^x(lax)
                =la z_0=lz=\Phi_z^y(ly)                \tag{C.51.8}
\]
identifies the actual maps, with the same target orbit.
The orbit is contained in $\partial O$, since $L$
preserves $O$ and its boundary. It is nonopen; an
open orbit meeting $\partial O$ would be an open
neighborhood disjoint from $O$, contradicting the
definition of its closure. Thus C.28 makes a
nontrivial such orbit an intrinsic line or circle.
For that case (C.51.8) is exactly one of the original
boundary quotients whose lifts define $\mathscr F_O$.
It follows in particular that an empty $\mathscr F_O$
forces every frontier orbit to be a point: apply
the construction to a sequence approaching any
prescribed frontier point.

\emph{The intrinsic family and every boundary equation.}
For $|\mathscr F_O|=1$ use the contraction $K^O_{s,t}$
of C.48, and for $|\mathscr F_O|=2$ use C.46.
Both commute with $H_O$, hence with its restricted
closure $Q$, and both have jointly continuous inverses.
They retain every quotient equation
\[
                    \Phi K^O_{s,t}=Q^B_{s,t}\Phi,
                                                        \tag{C.51.9}
\]
where $Q^B$ is the C.34 median family on the actual
boundary orbit $B$, with its based real covering
lift when that orbit is a circle.
For an empty family, C.32 makes $f_t|_O$ null in
$\operatorname{Homeo}(O)$ and C.33 supplies a based
null-homotopy $K^O$ in its exact $Q$-centralizer.
All frontier orbits are then points, as just proved,
so (C.51.9) for each constant $\Phi$ is automatic.

Define $\overline K^O$ to be $K^O$ on $O$ and
the C.34 family on $\partial O$. We prove its full
joint continuity, rather than continuity only for
fixed parameters. Let
\[
       x_n\in O,\quad x_n\to z\in\partial O,
       \qquad(s_n,t_n)\to(s,t).
                                                        \tag{C.51.10}
\]
Use $x_n=q_ny_n$ and a subsequence $y_n\to y$ as
above. The set
\[
 S=C\ \cup\!\bigcup_{(u,v)\in[0,1]^2}
                    \bigl(K^O_{u,v}(C)\cup(K^O_{u,v})^{-1}(C)\bigr)
                                                        \tag{C.51.11}
\]
is compact in $O$, by joint continuity on a compact
parameter space. Refine the choice of actual operators
in (C.51.4) to satisfy
\[
                       \sup_{w\in S}|h_nw-q_nw|<1/n.
                                                        \tag{C.51.12}
\]
For each $n$ this is a compact-open neighborhood
condition on $q_n\in Q$, so density of the specified
restrictions supplies $h_n\in H_O$.
The previous argument still proves (C.51.6)--(C.51.8).

Commutation with $Q$ and (C.51.12) give
\[
 \begin{split}
 K^O_{s_n,t_n}(x_n)&=q_nK^O_{s_n,t_n}(y_n),\\
 \bigl|K^O_{s_n,t_n}(x_n)
            -h_nK^O_{s_n,t_n}(y_n)\bigr|&<1/n.
 \end{split}                                            \tag{C.51.13}
\]
The points $K^O_{s_n,t_n}(y_n)$ converge in $O$ to
$K^O_{s,t}(y)$ and lie in the compact set $S$.
Uniform convergence (C.51.7), followed by the exact
quotient identity (C.51.9), therefore yields
\[
 \begin{split}
 h_nK^O_{s_n,t_n}(y_n)
      &\longrightarrow\Phi_z^y K^O_{s,t}(y)\\
      &=Q^{Gz}_{s,t}(\Phi_z^y(y))
        =Q^{Gz}_{s,t}(z).
 \end{split}                                            \tag{C.51.14}
\]
For a point orbit the same computation uses the
constant map and gives $z$.
Every subsequence in (C.51.10) has a further subsequence
of this form because $C$ is compact. The value on the
right is independent of all choices. Thus (C.51.14)
holds for the entire sequence and proves continuity
from $O$ to the prescribed frontier value.
Continuity for sequences contained in $\partial O$
is precisely C.34. Splitting a sequence according
to its two pieces proves joint continuity everywhere
on $[0,1]^2\times\overline O$.

The inverse maps have the same argument, using
the inverse images already included in $S$.
Equation (C.51.9) implies
$\Phi(K^O_{s,t})^{-1}=(Q^B_{s,t})^{-1}\Phi$
by composing with these actual inverse maps.
The inverse boundary family is jointly continuous
by C.34. This proves the inverse version of
(C.51.14), and hence joint continuity of the
piecewise inverse on the entire closure.
The maps preserve $O$ and each frontier orbit and
are bijections there; the piecewise inverses are
their actual inverses. They are homeomorphisms.
The endpoint and based-edge identities and
commutation with $H_O$ hold on both pieces.

\emph{Transport and finite pasting.}
The class of orbits in the claim is $H$-invariant:
conjugation by $h\in H$ identifies the two exact
groups $Q_O,Q_{hO}$, their quotient spaces, and
their full boundary-quotient families by C.46.
Choose one orbit from each $H$-orbit of these open
components, make the construction there, and set
\[
             \overline K^{hO}_{s,t}
                   =h\overline K^O_{s,t}h^{-1}.
                                                        \tag{C.51.15}
\]
The equality with every stabilizing operator makes
this independent of the choice of $h$ carrying the
selected representative to a given component.
Median naturality in C.27--C.34 makes its frontier
restriction the identical prescribed C.34 family.

Distinct open-orbit closures meet only in the
closed nonopen locus $Y$. Every one of the new
families and its inverse equals the corresponding
C.34 map there. The finite closed-set pasting
lemma thus gives the asserted maps and inverse
maps on $Y$ together with finitely many closures.
For every finite collection $\mathcal I$, the
homeomorphism $h$ carries this union to the union
indexed by $h\mathcal I$ and intertwines their
homotopies by (C.51.15). This is the exact
operator-transport comparison on those finite unions.

The proof uses one fixed compact fundamental set
inside each treated orbit. It supplies no uniform
control over infinitely many different such sets,
and it does not identify the median boundary value
with the exponential value forced on $X_*$ by C.50.
These are explicit remaining comparisons in C.12.
\end{proof}
